\documentclass[language=spanish]{../../meetingminutes}

\headerlogo{../../logo.pdf} % Path to header logo image
\institution{LinuxUMU}
\department{}

\title{Estatutos Asociativos}
\type{} % Built-in types: BoardMeeting, GeneralAssembly, ExtraordinaryGeneralAssembly, or custom string
\date{12 de Octubre, 2026}
\location{}
\starttime{}
\endtime{}

\begin{document}

\maketitle

\section*{CAPÍTULO I: DENOMINACIÓN, CONSTITUCIÓN, ÁMBITO, FINES, DURACIÓN, DOMICILIO.}

\textbf{Art. 1.-} Con la denominación (*) de LinuxUMU se constituye en una Asociación sin ánimo de lucro,
de ámbito UNIVERSITARIO y capacidad plena de obrar, de acuerdo con lo establecido en el artículo 22 de la Constitución,
la Ley Orgánica 1/2002, de 22 de marzo, reguladora del Derecho de Asociación y demás disposiciones vigentes que regulen la materia.

\textbf{Art. 2.-} Serán \textbf{fines} de la misma

\textbf{Art. 3.-} Para el cumplimiento de los fines del artículo anterior,
la entidad promoverá y organizará diversas \textbf{actividades} como:

Las actividades necesarias para el cumplimiento de sus fines, habrán de atenerse
a la legislación específica que regule tales actividades.

Los beneficios obtenidos por la Asociación, derivados del ejercicio de actividades
económicas, incluidas las prestaciones de servicios, deberán destinarse,
exclusivamente, al cumplimiento de sus fines, sin que quepa en ningún caso su
reparto entre los asociados ni entre sus cónyuges o personas que convivan con
aquellos, con análoga relación de afectividad, ni entre sus parientes, ni su
cesión gratuita a personas físicas o jurídicas con interés lucrativo.

\textbf{Art. 4.-} La duración de la entidad es por tiempo indefinido,
pudiendo ingresar en ella nuevos/as socios/as o causar baja los/las antiguos/as
sin necesidad de nueva constitución.

\textbf{Art. 5.-} El domicilio social se encuentra en: Sala Asociativa 11,
Centro Social Universitario (Edificio 29), Campus de Espinardo, 30100 Murcia

\section*{CAPÍTULO II: DE LOS/LAS SOCIOS/AS, DERECHOS Y DEBERES, PROCEDIMIENTOS, ADMISIÓN Y PÉRDIDA DE LA CUALIDAD DE SOCIO/A.}

\textbf{Art. 6.-} Pueden asociarse todas las personas mayores de edad y con capacidad de obrar que deseen cooperar con sus fines y acepten los presentes Estatutos, sin perjuicios de la admisión de los/las menores no emancipados/as de más de catorce años que con el consentimiento, documentalmente acreditado de las personas que deban suplir su capacidad, se integren en secciones juveniles cuyos miembros tendrán voz pero no voto, no pudiendo ser miembros de la junta directiva.
 La organización interna y el funcionamiento de la entidad será democrática, con pleno respeto al pluralismo. 

\textbf{Art. 7.-} Los/las socios/socias tienen los siguientes derechos:

7.1.- Participación en las actividades de la entidad y en los órganos de gobierno y representación, a ejercer el derecho de voto, así como asistir a la Asamblea General, de acuerdo con los Estatutos.
7.2.- A ser informados/as acerca de la composición de los órganos de gobierno y representación de la entidad, de su estado de cuentas y del desarrollo de su actividad.
7.3.- A ser oído/a con carácter previo a la adopción de medidas disciplinarias contra él/ella y a ser informado/a de los hechos que den lugar a tales medidas, debiendo ser motivado el acuerdo que, en su caso, imponga la sanción.
7.4.- A impugnar los acuerdos de los órganos de la entidad que estime contrarios a la ley o a los Estatutos.

\textbf{Art. 8.-} Los/las socios/as vienen obligados/as a:

8.1.- Compartir las finalidades de la entidad y colaborar para la consecución de las mismas.
8.2.- Pagar las cuotas, derramas y otras aportaciones que, con arreglo a los Estatutos, puedan corresponder a cada socio/a.
8.3.- Cumplir el resto de obligaciones que resulten de las disposiciones estatutarias.
8.4.- Acatar y cumplir los acuerdos válidamente adoptados por los órganos de gobierno y representación de la entidad.

\textbf{Art. 9.-} La condición de socio/a se pierde:

9.1.- Con actos que perjudiquen  gravemente los intereses de la entidad.
9.2.- Por propia voluntad.
 
En todo caso, será la Asamblea General quien decida sobre las posibles expulsiones de socios/as.

\textbf{Art. 10.-} Procedimiento de admisión de socios/as:

10.1.- La admisión de socios/as podrá ser decidida por la Junta Directiva y siempre que cumplan las condiciones previstas en los Estatutos.

10.2.- En todo caso, la Asamblea General ratificará o revocará las decisiones de la Junta Directiva.

10.3.- La entidad dispondrá  de una relación actualizada de las personas asociadas y recoger en un libro las actas de las reuniones de sus órganos de gobierno y representación. 




\end{document}
